\begin{lstlisting}[language=c++,label={code:27_03},caption={回帰分析をやってみる}]
# データを入力
koukou <- c(2.13,2.42,2.26,3.87,3.90,2.43,3.44,
2.15,2.18,3.00,3.42,2.55,3.19,3.05,2.52)
daigaku <- c(460,500,473,620,690,512,582,
550,485,650,593,528,585,569,518)
# データフレーム型に組み上げる
seiseki <- data.frame(koukou,daigaku)

# データは図にする
ggplot(seiseki,aes(x = koukou, y = daigaku)) +
    geom_point() +
    geom_smooth(method = "lm", se = F)

# 線形モデルはlm関数。従属変数~説明変数の形で書く
fit <- lm(daigaku ~ koukou, data = seiseki)
# 結果をわかりやすく表示
summary(fit)
\end{lstlisting}
